\documentclass{article}
\usepackage{graphicx} % Required for inserting images

\title{Rapport-Fonseca-Fourmond}
\author{Antoinette Fourmond, Gabriel Fonseca Rodrigues }
\date{February 2026}

\begin{document}

\maketitle

\section{Introduction}
Notre binôme développe un algorithme de réseau de neurones simplifié (supervisé puis semi-supervisé) pour détecter des anomalies et tentatives d'intrusion 
sur un réseau.

Nous avons choisi ce sujet pour apprendre les principes fondateurs de l’intelligence artificielle, qui prend de plus en plus de place dans notre quotidien.

En effet, ce projet a aussi pour but de détecter des patterns définis comme des anomalies par la base de données. On souhaite y ajouter un algorithme de 
prétraitement de données pour qu’il soit plus précis et/ou rapide.

Pour que ce projet soit considéré comme un succès, notre binôme a posé trois métriques d'évaluations : 
La première, considère l'efficacité intrinsèque de l'algorithme développé -  c'est-à-dire sa précision à répondre que les données d’entrées sont des anomalies
définies par la base. Nous considérons que sa valeur doit être supérieure à 95\% pour que ce projet soit réussi. Parallèlement, la perte sera la plus proche 
possible de zéro même si cette donnée n’est là que pour apporter de la cohérence et non pas une évaluation.

On s'intéresse à la contribution de notre projet pour la communauté scientifique, on souhaite donc comparer la précision de notre algorithme avec celle d'autres
algorithmes de détection déjà existants.


\section{Implémentation}
Nous avons réalisé notre projet en suivant la méthode SCRUM. En plus des réunions avec vous, les professeurs chaque semaine, on en avait aussi entre nous.
Chaque tâche de notre projet a été divisée en deux parties : le cadrage puis l’implémentation. Ces sprints nous permettent donc soit de rassembler toutes 
les idées de la conceptualisation au résultat final de la tâche, soit de réaliser cette dite tâche.

Nous avons codé en langage python, utilisant les librairies "numpy", pour les conversions et le calcul matriciel des données et "pandas", pour le 
traitement et nettoyage des données d'entraînement et tests et "sklearn.preprocessing", pour l'encodage et normalisation des données de type "Object".
Le projet est jusqu'à présent divisé sur deux fichier principaux, "main.py" et "protoN.py".

\subsection{Main}
Ce fichier contient les extractions et traitements des données pour être injectées en entrée sur le réseau de neurones. 

Notre fonction “prepare\_data”, récupère les données .csv pour les mettre dans un DataFrame. On encode les colonnes dont le type est Objet. On retire les
colonnes dont les données sont arbitraires malgré l’encodage, comme les adresses IP. On a décidé d'entraîner notre modèle sur 75\% des données pour éviter 
l’overfitting et le testons sur 25\%.

Il est aussi très important de normaliser l'ensemble des données avant l'exécution. On a d’ailleurs eu pendant longtemps des soucis, car on renormalisait à
chaque mini-bacth de données. Cela avait pour conséquence que les poids appris sur un batch deviennent incohérents sur le suivant. Le modèle ne pouvait pas 
converger.

"slice\_per\_time" partage les paquets de la base des données en fonction d’une fenêtre de temps changeante dynamiquement. L’objectif est de simuler à notre
algorithme un apprentissage par flux des données.


\subsection{ProtoN}

Notre projet a pour but de réaliser  un algorithme de réseau de neurones simplifié. La classe "Neuron", définit la structure d'un neurone avec ses fonctions
d'entraînement, d'évaluation, d'activation, de prédiction et d'optimisation. 

La première difficulté rencontrée est qu’on ne connaissait pas beaucoup de concepts de ce domaine. On a dû s'exercer sur des réseaux plus simples comme la 
classification de plantes toxiques pour voir comment les différents modules fonctionnent eux.
La fonction d'activation  appelée "sigmoid" donne la probabilité qu’une donnée d'entrée soit une anomalie entre [0-1]. Sa formule est définie comme :
\( 1/1+e^{-z} \), pour z la sommes des valeurs d'entrées par ses poids plus ses bias \(z = \Sigma(x_i*w_i)+b \).

La fonction d'optimisation, aussi appelée descente de gradient, a comme objectif d'améliorer les futurs calculs du réseau par l'ajustement des poids et du 
biais. L'ajustement des poids se f: \(w=w- lr*(1/length) * d * (m-y)\) et pour l'ajustement du biais:

\(b=b-lr*(1/length) * \Sigma_i m_i - y_i \),  où $lr$ est défini comme le taux d'apprentissage (learning rate). Le rôle de cette fonction est de minimiser
l'erreur globale en déplaçant les paramètres du neurone dans la direction opposée au gradient de la fonction de perte.

Un de nos bugs majeurs a été la prédiction. On comparait les probabilités de la sigmoid à un seuil calculé sur les pertes (loss) ce qui bloquait l'accuracy 
à 50\%. La correction a simplement comparé la sortie à 0.5 : si > 0.5 → anomalie, sinon → normal.

Après avoir implémenté nos fonctions principales et réglé nos bugs énoncés plus haut, la précision de notre modèle était à 72\% .

Nous avons choisi d'implémenter la méthode ADAM (Adaptive Moment Estimation) pour optimiser nos pas. ADAM combine les avantages de deux autres méthodes : 
elle utilise la moyenne des gradients précédents pour accélérer la descente (comme un élan), et la variance des gradients pour adapter le taux d'apprentissage
à chaque paramètre. Concrètement, cela permet au neurone d'avancer plus vite dans les directions stables et de ralentir là où les gradients oscillent, tout en
corrigeant le biais des premières itérations. Résultat : une convergence plus rapide et plus stable qu'une descente de gradient classique. Cela a augmenté la 
précision de notre modèle à 76.6\%.
Il est également important de diminuer au maximum le taux de faux négatifs. Sur l’ensemble des données, sans regarder les sous-ensembles (bach), on a 16\% de
faux positifs et 33\% de faux négatifs.


\section{Jalon}

Notre méthodologie SCRUM nous a permis de réaliser nos premiers objectifs dans les temps qu’on s’était donné. En effet, nous avions prévu la fin de l’implémentation
de notre réseau de neurones à une couche et de la méthode ADAM pour le 19 mars. Ce temps a servi à avoir un prototype solide qui sert de base à nos prochains ajouts.

Premièrement, nous avons pour objectif d’ici mi-mars, que notre modèle ait une précision supérieure à 95\%. La baisse de la perte viendra parallèlement mais elle
n’est pour nous pas un facteur d’évaluation.  

Un perceptron à un seul neurone ne peut tracer qu'une frontière de décision linéaire. Les patterns d'intrusion réseau impliquent des relations non-linéaires 
complexes entre débit, nombre de paquets, durée, etc. Pour cette raison, nous envisageons l’ajout d’une ou deux couches cachées avec une activation ReLU. Celle-ci
rendra notre modèle plus performant par sa capacité à détecter des motifs complexes, en transformant l’espace pour que les combinaisons d’entrées (motifs d'attaque)
deviennent linéairement séparables. 

Une fois cette première tâche réalisée, nous souhaitons comparer l’accuracy de notre modèle avec celle d’algorithmes de ML supervisés et non supervisés tel 
qu'Isolation Forest et K-Mean clustering. 
Parallèlement, nous sommes déjà en train de réfléchir à un algorithme à intégrer à notre réseau. Nous avons plusieurs pistes, comme le prétraitement de données. En 
effet, les IP étant des données définies arbitrairement, nous ne les prenons pas actuellement en compte. Notre première piste est de les traiter différemment pour 
qu’elles puissent apporter les informations qu’elles contiennent. 

Comme expliqué dans la partie implémentation, nous avons une fonction "slice\_per\_type" qui découpe les données par tranches de temps. En deuxième piste, on souhaiterait 
approfondir ce découpage dynamique pour qu’il découpe de manière plus intelligente. 

Une troisième piste serait l’introduction de la Markov Chain.

\end{document}
